\begin{scanpage}{0082}{65}
\[\mathbf A=\frac1c\sqrt{\frac{2\pi}{k_1r_1}}e^{ik(x_1\cos\alpha_1+r_1\sin\alpha_1)+i\pi/4}\int_{-\infty}^a\mathbf J(\eta)e^{ik\eta\Phi_1}\,d\eta\]
by a replacement of
\[\int_0^a\mathbf J(\rho,\psi)e^{ik\rho\Phi}\,d\rho\quad\text{by}\quad\int_{-\infty}^a\mathbf J(\eta)e^{ik\eta\Phi_1}\,d\eta\]
in those cases when $\Phi=\Phi_1$. The functions $\Phi$ and $\Phi_1$ are determined by the equations
\[
\left.\begin{aligned}
\Phi&=\sin\gamma\sin\psi-\sin\vartheta\cos(\psi-\varphi),\\
\Phi_1&=\sin\gamma\sin\psi-\sin\vartheta\cos(\psi-\varphi)\sqrt{\frac{1-\sin^2\gamma\cos^2\psi}{1-\sin^2\vartheta\sin^2(\psi-\varphi)}}.
\end{aligned}\right\}\tag{11.34}
\]
Solution (11.32) was determined exactly in this way with $ka\sqrt{\lambda^2+\mu^2}\gg1$, when, for auxiliary half-plane whose edge touches the rim of the disk at the points $\psi=\delta$, $\psi=\pi+\delta$, the phase $\Phi_1$ was equal to $\Phi$.

A solution to the problem using this method also is possible in the case
\[\vartheta=\gamma,\quad\varphi=\frac\pi2,\tag{11.35}\]
when $\Phi_1=\Phi=0$. The direction $\vartheta=\gamma$, $\varphi=\frac\pi2$ corresponds to the principal maximum of the scattering diagram, and therefore is of special interest. Substituting the relationships
\[
\left.\begin{aligned}
I_{y_1}&=\frac{c}{ik4\pi}\frac{H_{0x_1}}{\sin\alpha_1}\frac{\sin\varphi_1^0-1}{\cos^2\varphi_1^0},\\
I_{x_1}&=\frac{c}{ik4\pi}\frac1{\sin^2\alpha_1}(E_{0x_1}+\cos\alpha_1\operatorname{tg}\varphi_1^0 H_{0x_1})\frac{\sin\varphi_1^0-1}{\cos\varphi_1^0},
\end{aligned}\right\}\tag{11.36}
\]
which follow in this case from (11.19) into the equations
\[
\left.\begin{aligned}
A_x&=\frac ac\frac{e^{ikR}}R\int_0^{2\pi}(I_{y_1}\cos\psi+I_{x_1}\sin\psi)\,d\psi,\\
A_y&=\frac ac\frac{e^{ikR}}R\int_0^{2\pi}(I_{y_1}\sin\psi-I_{x_1}\cos\psi)\,d\psi,
\end{aligned}\right\}\tag{11.37}
\]
\end{scanpage}
